
\Universita{Università di Trento}
\Dipartimento{Dipartimento di Ingegneria e Scienze dell' Informazione}
\CorsoDiLaurea{Laurea Triennale in Informatica }
\AnnoAccademico{Anno Accademico 2025--2026}
\Titolo{\textbf{Quantum Annealing Learning Search per la risoluzione di problemi QUBO}}
\Relatore{Prof. Enrico \textsc{Blanzieri}}
\RelatoreLabel{Relatore}
\CandidatoLabel{Candidato}

\Candidato{Francesco \textsc{Buscardo}} 
\Matricola{246031}
\DataEsame{1 Novembre 2026}
\Logo{Impaginazione/logo_rosso.png}
\LogoWidth{4.5cm} %optional, default: 3cm
\LogoPosition{top}
\LogoSfondo{Impaginazione/logo_bn.png}
\opacitaSfondo{0.05}

\begin{titlepage}
    \newgeometry{left=3cm, right=3cm, bottom=2cm, top =3cm} 
    \pagestyle{empty}
    \makefrontpage
    \restoregeometry
\end{titlepage}


\newpage
\thispagestyle{empty}
\mbox{}  % This ensures an empty page
\newpage


\frontmatter
\null\vspace{\stretch{1}}
\begin{flushright}
    \textit{Dedica}
\end{flushright}
\vspace{\stretch{2}}\null

\chapter*{Ringraziamenti}

\begin{abstract}
L'obiettivo del lavoro è esplorare come Quantum Annealing Learning Search (QALS) riesca a risolvere problemi di ottimizzazione quadratica binaria (QUBO) e confrontare le sue prestazioni con altri metodi classici.
\end{abstract}

\tableofcontents

\printglossaries
\addcontentsline{toc}{chapter}{Glossario}

\printnomenclature
\addcontentsline{toc}{chapter}{Nomenclature list}

\mainmatter
\chapter*{Introduction}
\markboth{\MakeUppercase{Introduzione}}{}
\addcontentsline{toc}{chapter}{Introduzione}
\input{capitoli/intro}